\chapter{Estado del Arte}
\label{cap:estado_arte}


\section{Método Clásico}
\label{sec:metodo_clasico}
El análisis y construcción de carteras eficientes ha sido un área central en las finanzas 
cuantitativas desde la propuesta original de Markowitz\cite{markowitz1952}, donde se formalizó la 
optimización media–varianza como problema fundamental de selección de activos. A pesar de su 
relevancia histórica, numerosos autores han demostrado que estos modelos clásicos sufren 
importantes limitaciones en entornos reales: inestabilidad numérica, sensibilidad extrema al ruido 
en las estimaciones de la matriz de covarianza y falta de robustez en presencia de correlaciones 
cambiantes. López de Prado muestra en \cite{lopezdeprado2019frontier} que estos métodos suelen 
amplificar el error de estimación en la matriz de covarianza en lugar de mitigarlo, lo que provoca 
carteras inestables y poco realistas. Estas limitaciones motivan el desarrollo de enfoques 
alternativos como HRP.

\section{HRP Clásico y Variantes No-Cuánticas}
\label{sec:hrp_clas_y_variantes_no_cuanticas}
Estas limitaciones dieron lugar al desarrollo de alternativas más robustas como el Hierarchical 
Risk Parity (HRP)\cite{lopezdeprado2016hrp}. HRP evita la inversión de la matriz de covarianza y 
utiliza en cambio técnicas de clustering jerárquico para organizar los activos según distancias 
obtenidas a partir de la correlación. Una vez construida la estructura jerárquica, el riesgo se 
reparte mediante un procedimiento de asignación por bisección recursiva (recursive bisection) que 
genera carteras más estables y menos sensibles al ruido. 
No obstante, HRP sigue dependiendo de la definición de la métrica de distancia, lo que introduce 
una dependencia crítica en la representación de las relaciones entre activos.
Desde su presentación, numerosas extensiones se han propuesto, incluyendo variantes con 
restricciones\cite{pfitzinger2019constrainedhrp}, métricas de distancia alternativas 
\cite{salasmolina2025hrpdist} y métodos no lineales basados en kernels.

\section{KHRP}
\label{sec:khrp}
Una de estas extensiones es el Kernel HRP (KHRP), que sustituye la correlación estándar por 
distancias derivadas de embeddings no lineales mediante kernels (p. ej. RBF) y métricas inducidas 
como MMD. El uso de kernels permite capturar relaciones complejas entre activos que la correlación 
lineal no refleja adecuadamente\cite{sriperumbudur2011kernels}. Trabajos recientes han mostrado que 
la elección de la métrica de distancia influye de forma decisiva en la estructura jerárquica del 
HRP y en el rendimiento de la cartera\cite{salasmolina2025hrpdist}, lo que ha impulsado la búsqueda 
de representaciones más expresivas.

\section{Computación Cuántica en Finanzas}
\label{sec:computacion_cuantica_en_finanzas}
Paralelamente, la investigación en computación cuántica aplicada a las finanzas ha experimentado un 
crecimiento significativo en los últimos años. Los primeros trabajos, como \cite{rebentrost2018quantumportfolio}, 
demostraron que algoritmos cuánticos como Quantum Amplitude Estimation (QAE) permiten acelerar 
ciertos cálculos financieros, especialmente los basados en Monte Carlo. Egger y Woerner 
\cite{egger2019creditrisk} aplicaron estas técnicas al análisis de riesgo de crédito siguiendo el 
marco de Basilea II, estableciendo así una conexión directa entre computación cuántica y 
metodologías cuantitativas tradicionales. Otros avances, como Real-QAE \cite{manzano2023rqae} o el 
análisis del sesgo en QAE iterativa \cite{miyamoto2024iqae}, han contribuido a disponer de 
herramientas cuánticas más estables y adecuadas para hardware NISQ. Una revisión sistemática de 
este campo puede encontrarse en \cite{naik2023quantumfinance}. En general, estos trabajos se 
centran en tareas de estimación y valoración financiera, más que en la estructuración jerárquica 
de carteras del tipo HRP.
Sin embargo, la mayoría de estos enfoques se encuentran limitados por las restricciones del 
hardware NISQ, incluyendo ruido, baja fidelidad y número reducido de qubits.

\section{QHRP}
\label{sec:qhrp}
En este contexto surge el trabajo más reciente de Arian et al.\cite{arian2025qhrp}, Quantum Hierarchical Risk 
Parity (QHRP), que constituye el punto central del presente Trabajo de Fin de Grado. QHRP 
reinterpreta la etapa fundamental del HRP, la construcción de la matriz de distancias entre 
activos, a través de un enfoque cuántico basado en feature maps cuánticos y matrices de densidad 
promedio. En lugar de medir correlaciones o distancias clásicas, QHRP calcula distancias cuánticas 
obtenidas del grado de similitud entre estados cuánticos codificados a partir de los datos 
financieros. Este enfoque proporciona una métrica definida en el espacio de estados cuánticos (espacio de Hilbert), 
con mayor capacidad para capturar relaciones no lineales entre activos que sus alternativas clásicas. 
Según los experimentos reportados por los autores y en 
los escenarios considerados, QHRP supera a HRP clásico y KHRP en términos de estabilidad, Sharpe 
ratio y robustez, utilizando datos reales de mercado. Otro aspecto clave del trabajo es que 
proporciona un repositorio GitHub con código completamente reproducible, lo que facilita su 
implementación y la experimentación práctica.

\section{Herramientas Cuánticas Utilizadas}
\label{sec:herramientas_cuanticas_utilizadas}
Por último, dado que QHRP se basa en técnicas cuánticas inspiradas en el modelo gate-based, 
la herramienta de referencia para la construcción y simulación de circuitos cuánticos es Qiskit. 
Esta librería de IBM se ha convertido en el estándar de facto para la programación cuántica debido 
a su acceso directo a hardware real, su comunidad activa y la amplia documentación disponible. En 
particular, el ecosistema incluye simuladores como Aer, que resultan fundamentales para validar y 
comparar circuitos en condiciones controladas antes de su ejecución en dispositivos físicos. 
Qiskit, además, incluye módulos específicos para finanzas que implementan algoritmos como QAE o 
modelos de riesgo crediticio, constituyendo un entorno adecuado para familiarizarse con los 
conceptos necesarios para comprender QHRP en el contexto actual de hardware NISQ. Por este motivo, 
y siguiendo la práctica habitual en la comunidad, se utilizará Qiskit como entorno de simulación y 
experimentación cuántica en este proyecto.

